Apart from the electron-electron [\cref{sect:GW_approximation}] and phonon-phonon [\cref{sec:anharmonic_vibs}] interactions, the electron-phonon interactions (EPI) give rise to another many-body quantum effect that plays an important role in dictating the physical properties of materials. The fundamental concept of electronic band structures [\cref{sect:band_structure}] is formulated based on the Bloch's theorem which considers the electronic wavefunction in a solid is modulated by the ionic potential $V(\bm{R})$ that is perfectly periodic in space, such that electrons moving in this potential should experience no scattering from the lattice. It is apparent that such an idealisation breaks down at finite temperature, as the phonon vibrations in the crystal will deform the crystal potential, as $\Delta V(\bm{R}_i)\sim \Delta \bm{R}_i\nabla_{\bm{R}}V(\bm{R}_{i})$ to the first-order approximation, which contributes towards the coupling of electrons with acoustic phonons via the deformation potential.\cite{khan1984deformation} Particularly, in noncentrosymmetric crystals where the long-range cooperative acoustic phonon vibrations can also generate electric fields due to the piezoelectric effect that couples with the electrons, leading to the piezoelectric EPIs.\cite{zook1964piezoelectric}   Finally, the oscillating electric dipole moments induced by the optical phonon vibrations in polar crystals can also couple with electrons, leading to the so-called Fr\"{o}hlich interaction\cite{frohlich1954electrons} being another important mechanism behind EPI in materials. 

Formally, the Hedin-Baym equations,\cite{hedin1965new,baym1961field} which are a set of integro-differential equations derived from the many-body quantum field theory, provided the most sophisticated description of the EPI effect in solids to-date, with harmonic phonons being the only assumption.\cite{giustino2017electron} In a high-level view, the Hedin-Baym equations take into account the  mutual effects of EPIs which alter the dynamics of both the electronic and phononic subsystems via the electronic ($\Sigma_{\mbox{\scriptsize{eph}}}$) and phononic ($\Pi_{\mbox{\scriptsize{eph}}}$) self-energies, whereby the fluctuations in the lattice potentials introduce additional Coulomb screening on the electrons, which in terms also modifies the phonon  frequencies and lifetimes by changing the nuclear charge densities and nuclear-nuclear interactions. 

The complexity of the coupled integro-differential equations prevents the Hedin-Baym equations from being fully solved to obtain the self-consistent solutions for the energy shifts and lifetimes of the electronic and phononic subsystems with the presence of EPI. Nevertheless, the key physical quantity that determines the self-energies $\Sigma_{\mbox{\scriptsize{eph}}}$ and $\Pi_{\mbox{\scriptsize{eph}}}$ is the frequency- and momentum-dependent EPI strength $g_{mn\nu}(\bm{k},\bm{q})$, which couples the electronic Bloch states $|\psi_{m}(\bm{k}+\bm{q})\rangle$ and $|\psi_{n}(\bm{k})\rangle$ through a phonon vibration mode $\bm{\mathcal{Q}}_{\bm{q},\nu}$, where $m$, $n$ and $\nu$ are the band indicies and $\bm{k}$, $\bm{q}$ are the electronic and phonic wavevectors, respectively. The increasing computational power provided by the modern HPC facilities has made the computing $g_{mn\nu}(\bm{k},\bm{q})$ become readily accessible with major solid-state DFT codes,\cite{noffsinger2010epw,ponce2016epw,karsai2018electron} enabling us to gain critical insights into this many-body quantum effects through computational modellings. Here, we shall briefly account for the principles behind the two major state-of-the-art approaches to compute the EPI strengths in solids, which  follow directly  the real and reciprocal space formalism for harmonic phonon calculations as outlined in \cref{sect:phonon}. 

In the \emph{perturbative approach} based on DFPT,\cite{giustino2017electron} one works directly with the definition of the EPI matrix element
\begin{equation}
    \label{eq:epi_matrix_element_dfpt}
    g_{mn\nu}(\bm{k},\bm{q})=\left\langle\psi_{m}^{\mbox{\scriptsize{KS}}}(\bm{k}+\bm{q})\Bigg\vert \frac{\partial V^{\mbox{\scriptsize{KS}}}}{\partial \bm u_{\nu,\bm{q}}} \Bigg\vert  \psi_{n}^{\mbox{\scriptsize{KS}}}(\bm{k})\right\rangle,
\end{equation}
in which the computation of the perturbed KS potential $\partial V^{\mbox{\scriptsize{KS}}}$ follows the discussions in \cref{sect:phonon}. As discussed previously, DFPT would enable the direct computation of EPI strengths in which the contributions from long-wavelength ($\bm{q}\to0$) phonons are significant, without the need to use a large simulation supercell, as in the real-space \textit{finite difference} approach.\cite{monserrat2018electron} However, the later provides the additional flexibility whereby the lattice dynamics and electronic structures may be computed with two different levels of theories. In its essence, one recast \cref{eq:epi_matrix_element_dfpt} into a finite-difference form 
\begin{equation}
    \label{eq:epi_matrix_element_finite_diff}
    g_{mn\nu}(\bm{k},\bm{q})=\left\langle\psi_{m}^{\mbox{\scriptsize{KS}}}(\bm{k}+\bm{q})\Bigg\vert \frac{ V^{\mbox{\scriptsize{KS}}}(\Delta \bm u_{\nu,\bm{q}}) -V^{\mbox{\scriptsize{KS}}} (0) }{\Delta \bm u_{\nu,\bm{q}}} \Bigg\vert  \psi_{n}^{\mbox{\scriptsize{KS}}}(\bm{k})\right\rangle,
\end{equation}
which suggests that, under this formulation, one is required to perform separated calculations of the KS eigenstates at both the equilibrium structure (labelled as 0) and a series of  configurations that are generated by perturbing the equilibrium structure along different vibrational eigenstates $\{\bm{u}_{\nu,\bm{q}}\}$.

This approach tights up naturally with the core concept of PES sampling that underpins the current review. As pointed out by \citet{monserrat2018electron}, if one was interested in the change of some electronic property $\mathcal{A}$ (such as the band energy levels $\varepsilon_{n,\bm{k}}$ and band gap $E_g$) induced by the phonon vibration at temperature $T$ in general, without  consideration of the contributions from the specific vibrational modes or electronic states, we can first express $\mathcal{A}$ as a Taylor expansion of the vibrational coordinates, in the spirit of \cref{eq:higher_order_vib_Hamiltonian}:
\begin{equation}
    \label{eq:property_expand_vib_Hamiltonian}
    \mathcal{A}(\bm{R})=\mathcal{A}(0)+\sum_{\bm{q},\nu}\frac{\partial \mathcal{A}(0)}{\partial\bm{u}_{\bm{q},\nu}}\bm{u}_{\bm{q},\nu}+\sum_{\bm{q},\nu}\sum_{\bm{q}^\prime,\nu^\prime}\frac{\partial^2\mathcal{A}(0)}{\partial\bm{u}_{\bm{q},\nu}\partial\bm{u}_{\bm{q}^\prime,\nu^\prime}}\bm{u}_{\bm{q},\nu}\bm{u}_{\bm{q}^\prime,\nu^\prime}+\mathcal{O}(\bm{u}^3),
\end{equation}
which reflects how EPI arises from both harmonic and anharmonic phonons change the electronic behaviours of a material. If one truncates \cref{eq:property_expand_vib_Hamiltonian} at the harmonic term and taking the thermal average [in line with \cref{eq:general_sampling_2}], then, for example, the temperature-dependent shifts in the electronic state eigenenergies due to EPI at temperature $T$ will take the following form of Bose-Einstein averaging, leading to the famous Allen-Heine-Cardona (AHC) theory:\cite{allen1976theory,allen1981theory}
\begin{equation}
    \label{eq:electronic_state_epi_thermal_av}
    \varepsilon_{n\bm{k}}(T)=\varepsilon_{n\bm{k}}(0)+\sum_{\bm{q},\nu}\frac{1}{\omega_{\bm{q},\nu}}\frac{\partial^2 \varepsilon_{n\bm{k}}}{\partial \bm{u}^2_{\bm{q},\nu}}\left[\frac{1}{2}+\frac{1}{\exp(\hbar\omega_{\bm{q},\nu}/k_B T)-1} \right]
\end{equation}
in which $\varepsilon_{n\bm{k}}(0)$ is the electronic energy levels computed at 0 K. This provides an intuitive and semi-quantitative measure of the strength of EPI in semiconductors, as the shift of $\varepsilon_{n\bm{k}}$ (or equivalently, $E_g$) is directly proportional to $\frac{\partial^2 \varepsilon_{n\bm{k}}}{\partial \bm{u}^2_{\bm{q},\nu}}$. In addition to the physical origins of EPI as discussed above, the temperature-dependent variations in the electronic energy levels and band gaps can be understood as the results of modulating the orbital hybridisations between the bonding atoms in the solids due to the changes in the interatomic distances and angles from phonon vibrations. 

Another advantages of adopting the real-space approach is the capability to compute the anharmonic contributions to the EPIs by directly working with \cref{eq:property_expand_vib_Hamiltonian}. For example, \citet{saidi2016temperature} shows that the AHC theory based on harmonic phonon is insufficient to correctly capture the temperature-dependency of $E_{g}$ for  methylammonium lead iodide perovskite due to its strong vibrational anharmonicities. This can be achieved by performing AIMD simulations under the $NPT$ ensemble, one can access to a collection of time-dependent MD trajectories (frames) $\{\bm{R}(t)\}_T$ that correspond to the  configurations accessible on the fully anharmonic PES at a the target temperature $T$. By further computing the corresponding band gap energies on these MD frames, it becomes possible to retrieve the time-dependent evolutions of band gap energies $\{E_{g}(t)\}_T$ at the target temperature. From here, the expectation value of the band gap energies $\langle E_g(t)\rangle_T$ can be considered as the anharmonic equivalent to \cref{eq:electronic_state_epi_thermal_av}, providing information that is analogous to the real part of $\Sigma_{\mbox{\scriptsize{eph}}}$. At the same time, by computing the band gap dephasing function from $\{E_{g}(t)\}_T$ would further enable us to extract the band gap dephasing time due to the presence of EPI, which can be considered as a proxy for the imaginary part of $\Sigma_{\mbox{\scriptsize{eph}}}$.\cite{yang2020correlating,li2018time,li2019anharmonicity,hwang2004analysis}  

Apart from contributing to the temperature-dependent variation in $E_g$, EPI manifests itself in a wide range of other physical properties of materials, for which computational approaches is playing a significant role in facilitating the investigations of these phenomena. One of the most well-known effects is the phonon-assisted interband optical absorptions in the indirect band gap semiconductors, with Si being the most representative case,\cite{noffsinger2012phonon} which forms the basis of modern microelectronics and photovoltaics. In polar semiconductors, such as ZnO,\cite{teke2004excitonic} the photoexcited electron-hole pair can interact strongly via Coulomb interactions, forming quasiparticle called exciton, which can couple strongly to the longitudinal optical phonons that menifest in the appearance of multiple phonon sidebands in the optical spectra.

By solving Boltzmann transport equation [\cref{eq:e_cond_tensor}], taking into account the modification of the electron relaxation time due to EPI ($1/\tau\sim |g|^2)$, in which the EPI matrix elements $g$ was determined based on the Fermi-Golden rule, \citet{li2015electrical} demonstrated good agreement between the theoretically calculated and experimentally determined electron mobilities for semiconducting Si, metallic Al and 2D MoS$_2$, which can be considered a seminal work that providing a universal framework to include EPI in transport modellings for materials.  \citet{borysenko2010first} applied the DFPT scheme to compute the EPI matrix elements for graphene, which was further combined with the Boltzmann transport equation to demonstrate that both the acoustic and optical phonons contributed to the scattering of electrons at low temperature, whereas at the elevated temperature, the scatterings become dominated by the high-energy optical phonons. \citet{himmetoglu2014first} examined the strong contributions from optical phonon scattering to the electrical conductivities in \ce{SrTiO3}. Lately, \citet{mirjolet2021electron} provided the theoretical analysis on the $\rho(T)\sim  T^2$ dependency of resistivity for the strongly correlated perovskite of \ce{SrVO3}. Their investigation revealed that, in addition to the strong electron-electron coupling that renormalises the electron effective mass, the electronic band structure exhibits a cylindrical Fermi surface, in which a large portion of the Fermi surface can couple to a single phonon mode. This gives rise to polaron transport in the material due to strong EPI. 

EPI also plays the pivotal role in understanding the origins and predicting future candidate materials of high-temperature,  phonon-mediated superconductivity. It represents a type of superconductivity in materials (such as \ce{MgB2}\cite{joas2002theory} and FeSe/\ce{SrTiO3} interface\cite{yang2024phonon}), whereby the formation of electron pairs is mediated by the attractive polarisation fields that are generated by the lattice vibrations. The Migdal-Eliashberg theory,\cite{Migdal1958,Eliashberg1960} which extended the Bardeen-Cooper-Schrieffer (BCS) theory of super-conductivity\cite{Bardeen1957a,Bardeen1957b} that accounts for strong EPI. Upon the formation of electron pairing that contributes to the superconductivity in metals, its electronic DOS opens up a gap $(\Delta_{\mbox{\scriptsize{SC}}}$, commonly known as the superconducting gap. As the temperature rises beyond $T_c$, it closes the $\Delta_{\mbox{\scriptsize{SC}}}$, which destroys the electron pairing and the material returns back to its normal non-superconducting state. Based on this theory, McMilliam and Allan\cite{mcmillan1968transition,allen1975transition} established the following formula:
\begin{equation}
    \label{eq:mcmillan_Tc}
    k_B T_c=\frac{\hbar\omega_{\mbox{\scriptsize{log}}}}{1.2}\exp\left(-\frac{1.04(1+\lambda)}{\lambda-\mu^* (1+0.62\lambda)}\right),
\end{equation}
which provided the basic computational framework for determining the superconducting critical temperature $T_c$ for phonon-mediated superconductivity. Here, $\lambda$ is a material-specific parameter that is dependent on the frequency and momentum-dependent electron-phonon coupling strength, which can be determined from the isotropic Eliashberg function:\cite{allen1972neutron}
\begin{equation}
    \label{eq:eliashberg_function}
    \alpha^2 F(\omega)=\frac{1}{N_F}\int\frac{\mbox{d}\bm{k}\mbox{d}\bm{q}}{\Omega_{\mbox{\scriptsize{BZ}}}^2}\sum_{mn\nu}|g_{mn\nu}(\bm{k},\bm{q})|^2\delta(\varepsilon_{n\bm{k}}-\varepsilon_F)\delta(\varepsilon_{m,\bm{k}+\bm{q}}-\varepsilon_F)\delta(\omega-\omega_{q\nu}).
\end{equation}
It describes the modification of the electron and phonon density-of-states due to the presence of strong EPI in the materials system. From here, $\lambda=2\int_0^\infty\frac{\alpha^2 F(\omega)}{\omega}d\omega$, and the weighted logarithmic phonon frequency $\omega_{\mbox{\scriptsize{log}}}=\exp\left[\frac{2}{\lambda} \int_0^\infty\frac{\alpha^2 F(\omega)}{\omega}\log\omega d\omega\right]$. It should be immediately clear from \cref{eq:eliashberg_function} why strong EPI is prone to be observed in systems with cylindrical Fermi surfaces. Theoretically speaking, all ingredients to compute the spectral function with \cref{eq:eliashberg_function} are accessible with modern electronic structure theory. Practically, however, performing the integration would require sampling dense sets of electron and phonon wavevectors in the Brillouin zone, which can be seen as the main bottleneck in applying this method for studying materials with large and complex crystal structures.

Finally, the effect of EPI on the behaviour of the phonon subsystem should not be overlooked.\cite{giustino2017electron} Fundamentally, EPI contributes to the (first-order) phonon self-energy as
\begin{equation}
    \label{eq:phonon_self_energy}
    \Pi_{\bm{q}\nu}(\omega)=\frac{2}{N_{k}}\sum_{nm\bm{k}}\frac{|g_{nm\nu}(\bm{k},\bm{q})|^2[f(\varepsilon_{n\bm{k}})-f(\varepsilon_{m,\bm{k}+\bm{q}})]}{\hbar\omega+\varepsilon_{n\bm{k}}-\varepsilon_{m,\bm{k}+\bm{q}}+i\eta},
\end{equation}
in which $f(\cdot)$ is the Fermi-Dirac distribution function and $\eta\to 0^+$. From here, the shift in the phonon frequency from EPI can be computed as $\tilde{\omega}_{\bm{q}\nu}^2=\omega_{\bm{q}\nu}^2+2\omega_{\bm{q}\nu}\mbox{Re}\Pi_{\bm{q}\nu}(\tilde{\omega}_{\bm{q}\nu})$, and finite phonon lifetime $\gamma_{\bm{q}\nu}=-2\mbox{Im}\Pi_{\bm{q}\nu}(\tilde{\omega}_{\bm{q}\nu})$. It can be appreciated that when the effect of EPI on the phonon subsystem is significant, it can also lead to notable modifications on the lattice thermal conductivities of materials. A notable example of the effect of EPI on the phonon subsystem is the Kohn anomaly,\cite{kohn1959image} reflected as a sharp softening of phonon frequency. This effect is considered to be strong in 2D materials which exhibit a cylindrical Fermi surface and strong divergence of dielectric screening at the nesting wavevectors.\cite{zhu2015classification} For example, \citet{piscanec2004kohn} provided the theoretical analysis on the Kohn anomalies that are observed in graphite, in which they computed the phonon frequencies with DFPT at different levels of Fermi smearing, which affects the vertical transition probabilities between the occupied and virtual states, that can be considered as a manifestation of the EPI effect on the electronic subsystem.
